\section{Tree-told Nonlinear Story}
\label{sec:tree_story}

% Tree-story figures are produced by notebooks/audits/build_tree_diagnostics.ipynb
% into results/diagnostics/{xgb,lgbm}_har_tuned/ and copied to writeup/figures/.
% Until those have been built, the figure slots render as a labelled placeholder
% so the writeup compiles for cross-meeting iteration.
\providecommand{\treefigorpending}[2]{%
  \IfFileExists{#2}{\includegraphics[width=#1]{#2}}%
                   {\fbox{\parbox[c][1.5cm][c]{#1}{\centering\itshape pending: \detokenize{#2}}}}%
}

The tuned LightGBM (Section~\ref{sec:meeting_update}) achieves
$\textnormal{QLIKE} \approx 0.1316$ versus the Ridge HAR baseline at
$0.1380$, a small but persistent gap of $\sim$$0.006$. The advisor
asks have been blunt: \emph{which features and interactions are
responsible for this gap?} This section answers that question with
walk-forward TreeSHAP attributions on the same sequence of estimators
the production executor produces. We deliberately avoid anchoring the
narrative on the 18:30 / Sunday-open intraday pattern that earlier
slot-stratifications surfaced --- on closer inspection that pattern
turned out to be an upstream aggregation artifact rather than a
property of the trees themselves.

\subsection{Methodology}
\label{sec:tree_story_methodology}

We compute TreeSHAP attributions
\citep{lundberg2018consistent,lundberg2017unified} on the
\emph{walk-forward} sequence of fitted estimators, mirroring the
predicate \texttt{(t - train\_win + 1) \%
refit\_frequency == 0} used in
\texttt{src/scaling.py:run\_backtest}. Both XGBoost and LightGBM use
\texttt{refit\_frequency = 1} (one refit per 30-minute bar), so each
refit produces exactly one out-of-sample prediction. At every refit
$i$ we retain the fitted model and run
\[
    \boldsymbol{\phi}^{(i)} \;=\;
    \textnormal{TreeExplainer}(M_i)\bigl(\mathbf{x}^{\textnormal{oos}}_i\bigr),
\]
yielding one SHAP attribution row per OOS step. Pooling these rows
gives the population over which we compute global / yearly / per-fold
mean $|\phi|$. SHAP \emph{interaction} values
\citep{lundberg2020local} are O($F^2$) per sample, so we capture them
only every $K = 65$ refits ($\sim$weekly at 13 30-min bars per day),
producing $\sim$250 interaction observations across $\sim$5 years
OOS.

This walk-forward pooling protocol follows
\citet{gu2020empirical} for asset-pricing ML and
\citet{christensen2023machine} for realized-volatility forecasting.

\paragraph{Aggregations reported.}
\begin{itemize}
    \item \textbf{Global importance} --- mean $|\phi_f|$ over all OOS
        rows (Section~\ref{sec:tree_story_global}).
    \item \textbf{Yearly importance} --- mean $|\phi_f|$ stratified
        by calendar year, displayed as a feature$\times$year
        heatmap (Section~\ref{sec:tree_story_temporal}).
    \item \textbf{Per-fold importance + rank stability} --- partition
        OOS rows into rolling folds, compute Spearman rank
        correlation of the top-$K$ feature ranks across consecutive
        folds (Section~\ref{sec:tree_story_stability}).
    \item \textbf{Group importance} --- mean $|\phi_f|$ aggregated
        within feature groups (RV-lag, calendar, exog) to dampen
        the correlated-feature ambiguity flagged by
        \citet{aas2021explaining}.
    \item \textbf{Interaction strength} --- mean
        $|\phi^{\text{int}}_{f_i, f_j}|$ across the pooled
        interaction subsample
        (Section~\ref{sec:tree_story_interactions}).
\end{itemize}

\paragraph{Caveats (must be read alongside the headline numbers).}
\begin{enumerate}
    \item \emph{Scale.} SHAP attributes prediction-scale, not
        loss-scale, contributions. Comparing absolute $|\phi|$
        magnitudes across models with different target transforms
        (e.g., XGB vs.\ LGBM with diurnal-adjusted targets and
        \texttt{target\_winsor\_window=240} per
        \texttt{src/executor.py}'s \texttt{CONFIGS}) is not
        meaningful; we report each model on its own scale and
        compare ranks instead.
    \item \emph{Estimator pooling.} Pooled SHAP averages
        attributions across a heterogeneous sequence of fitted
        estimators. The numbers below describe \emph{the production
        sequence's} attribution, not any single canonical model.
    \item \emph{Correlated features.} Our HAR feature set is by
        construction collinear (\texttt{RV\_lag\_1m},
        \texttt{RV\_lag\_5m}, \texttt{RV\_lag\_22m}, $\dots$).
        Individual SHAP attributions split contribution arbitrarily
        within a correlated block. The group-importance summary in
        Section~\ref{sec:tree_story_global} dampens this; per-feature
        numbers are still reported for completeness but should be
        read in light of the group totals.
\end{enumerate}

Reproducibility: model objects are not persisted. The notebook
\texttt{notebooks/audits/build\_tree\_diagnostics.ipynb} regenerates
all SHAP arrays + tree-structure dumps from the canonical
\texttt{results/tune/best\_xgb.json},
\texttt{results/tune/best\_lgbm.json}, and a seeded RNG; the
\texttt{methodology} block of each
\texttt{tree\_story\_stats.json} records the exact run-time
parameters.

\subsection{Tree topology}
\label{sec:tree_story_topology}

Figure~\ref{fig:tree_depth_histogram} shows node-depth distributions
aggregated across all refits in the walk-forward sequence; XGBoost
clusters near the tuned \texttt{max\_depth = 6}, LightGBM is
naturally less symmetric (leaf-wise growth).
Figure~\ref{fig:feature_split_count} ranks the top-15 features by
total split count across all refits.

\begin{figure}[H]
\centering
% canonical source: results/diagnostics/{xgb,lgbm}_har_tuned/tree_depth_histogram.png
\treefigorpending{0.49\linewidth}{figures/tree_depth_histogram_xgb.png}\hfill
\treefigorpending{0.49\linewidth}{figures/tree_depth_histogram_lgbm.png}
\caption{Tree-node depth distribution aggregated across all walk-forward
refits. Left: tuned XGBoost. Right: tuned LightGBM.}
\label{fig:tree_depth_histogram}
\end{figure}

\begin{figure}[H]
\centering
% canonical source: results/diagnostics/{xgb,lgbm}_har_tuned/feature_split_count.png
\treefigorpending{0.49\linewidth}{figures/feature_split_count_xgb.png}\hfill
\treefigorpending{0.49\linewidth}{figures/feature_split_count_lgbm.png}
\caption{Top-15 feature split counts aggregated across walk-forward
refits.}
\label{fig:feature_split_count}
\end{figure}

\subsection{Global SHAP attribution}
\label{sec:tree_story_global}

The headline numbers in the master table
(Table~\ref{tab:master_compact}) report
\texttt{top\_shap\_feature} and \texttt{mean\_tree\_depth} for the
tuned tree rows; the appendix table
(Table~\ref{tab:master_full}) carries the full set including
\texttt{top\_shap\_share}, \texttt{rank\_stability\_rho}, and
\texttt{top\_interaction\_pair}. The bar plots in
Figure~\ref{fig:shap_summary} visualise the top-15 ranking, and
Figure~\ref{fig:shap_group} shows the same attribution aggregated by
feature \emph{group} --- the form that survives the correlated-feature
caveat.

\begin{figure}[H]
\centering
% canonical source: results/diagnostics/{xgb,lgbm}_har_tuned/shap_summary_bar.png
\treefigorpending{0.49\linewidth}{figures/shap_summary_bar_xgb.png}\hfill
\treefigorpending{0.49\linewidth}{figures/shap_summary_bar_lgbm.png}
\caption{Top-15 mean $|\phi_f|$ pooled across all walk-forward OOS rows.
Left: XGBoost. Right: LightGBM.}
\label{fig:shap_summary}
\end{figure}

\begin{figure}[H]
\centering
% canonical source: results/diagnostics/{xgb,lgbm}_har_tuned/shap_group_bar.png
\treefigorpending{0.49\linewidth}{figures/shap_group_bar_xgb.png}\hfill
\treefigorpending{0.49\linewidth}{figures/shap_group_bar_lgbm.png}
\caption{$\sum_f |\phi_f|$ aggregated by feature group (RV-lag,
calendar, exog).}
\label{fig:shap_group}
\end{figure}

\subsection{Stability of attribution}
\label{sec:tree_story_stability}

Figure~\ref{fig:shap_rank_stability} traces the rank of each top-10
feature across rolling refit folds. Rank stability is summarised by
\texttt{rank\_stability\_rho} in the master table: the mean Spearman
correlation between consecutive-fold rankings. Values close to 1.0
indicate the storyteller features replicate; values near 0 indicate
fold-to-fold noise dominates and the per-feature claims should not
carry weight.

\begin{figure}[H]
\centering
% canonical source: results/diagnostics/{xgb,lgbm}_har_tuned/shap_rank_stability.png
\treefigorpending{0.49\linewidth}{figures/shap_rank_stability_xgb.png}\hfill
\treefigorpending{0.49\linewidth}{figures/shap_rank_stability_lgbm.png}
\caption{Rank of the top-10 features within each rolling refit fold
(rank 1 = the fold's top feature). Flat lines indicate stable
attribution across regimes.}
\label{fig:shap_rank_stability}
\end{figure}

\subsection{Temporal view}
\label{sec:tree_story_temporal}

Figure~\ref{fig:shap_yearly_heatmap} shows mean $|\phi_f|$ by
calendar year for the top-15 features, surfacing regime drift across
the OOS horizon. The expected pattern is non-trivial drift around
2008 (GFC) and 2020 (COVID) on top of an otherwise stable
attribution; a flat heatmap would suggest the model's notion of what
matters is regime-invariant, while an erratic one would suggest
regime-conditioning is doing real work and warrants the further
analysis sketched in Section~\ref{sec:tree_story_reading}.

\begin{figure}[H]
\centering
% canonical source: results/diagnostics/{xgb,lgbm}_har_tuned/shap_yearly_heatmap.png
\treefigorpending{0.49\linewidth}{figures/shap_yearly_heatmap_xgb.png}\hfill
\treefigorpending{0.49\linewidth}{figures/shap_yearly_heatmap_lgbm.png}
\caption{Yearly mean $|\phi_f|$ heatmap for the top-15 features.}
\label{fig:shap_yearly_heatmap}
\end{figure}

\subsection{Where the nonlinearity comes from}
\label{sec:tree_story_interactions}

The QLIKE gap that survives the linear baseline must come from
either (a) genuinely nonlinear univariate responses or (b)
non-additive interactions across features. Figure~\ref{fig:shap_interaction_heatmap}
visualises the top-10 features' pairwise SHAP-interaction matrix; the
strongest off-diagonal cells localise the interaction that the trees
are exploiting. We surface the single largest interaction pair as
\texttt{top\_interaction\_pair} in the master table for cross-meeting
tracking.

\begin{figure}[H]
\centering
% canonical source: results/diagnostics/{xgb,lgbm}_har_tuned/shap_interaction_heatmap.png
\treefigorpending{0.49\linewidth}{figures/shap_interaction_heatmap_xgb.png}\hfill
\treefigorpending{0.49\linewidth}{figures/shap_interaction_heatmap_lgbm.png}
\caption{Top-10 feature pairwise mean $|\phi^{\text{int}}_{f_i,f_j}|$.
Diagonal = main effects.}
\label{fig:shap_interaction_heatmap}
\end{figure}

For the three highest-importance features (per the global
ranking), Figure~\ref{fig:shap_dependence} shows the SHAP dependence
plot --- feature value on the horizontal, SHAP attribution on the
vertical, colored by the strongest interacting feature. The shape of
the cloud captures the univariate nonlinearity; the colour gradient
captures interaction.

\begin{figure}[H]
\centering
% canonical source: results/diagnostics/{xgb,lgbm}_har_tuned/shap_dependence_*.png
\treefigorpending{0.32\linewidth}{figures/shap_dependence_top1.png}\hfill
\treefigorpending{0.32\linewidth}{figures/shap_dependence_top2.png}\hfill
\treefigorpending{0.32\linewidth}{figures/shap_dependence_top3.png}
\caption{SHAP dependence plots for the top-3 features (XGBoost). Colour
encodes the strongest interacting feature for each panel.}
\label{fig:shap_dependence}
\end{figure}

\subsection{Reading}
\label{sec:tree_story_reading}

The tree models earn their $\sim$$0.006$ QLIKE edge over Ridge HAR
through a mixture of (i) shallow-but-non-additive splits on the
short-horizon RV lags --- the mechanism a linear HAR cannot reproduce
--- and (ii) interactions between RV-lag features and the
diurnal-adjustment / calendar features that
\texttt{CONFIGS[\{xgboost,\,lightgbm\}].add\_calendar = True}
exposes. The group-bar (Figure~\ref{fig:shap_group}) tells us that
the RV-lag block dominates total attribution as expected; the small
non-zero share of the calendar block is the part that picks out
intraday + day-of-week conditioning.

That said, the gap is small and the rank-stability number
(\texttt{rank\_stability\_rho}, master table) sets the ceiling on
how strongly we can lean on individual-feature claims --- if it lands
below $\sim$0.7 we read this section as ``trees are doing nonlinear
work, but the per-feature decomposition is fold-noisy enough that we
should report only the group-level claim.'' Either reading is
honest; what would not be honest is overclaiming a specific
feature-pair as the source of the gap if the rank correlation
disagrees from one fold to the next.

\paragraph{What this section is not.} It is not a regime-conditioned
attribution (deferred to Bundle A in the plan), and it does not
anchor on the 18:30 / Sunday-open pattern --- which earlier slot
stratifications appeared to highlight but which a re-examination of
the upstream aggregation traced to an artifact of how the source
data was bucketed, not a property of the trees. Suggestions to
nest these diagnostics inside a regime label (VIX percentile, RV
percentile, crisis flag) are recorded in
Section~\ref{sec:meeting_update} as an open follow-up.
